\begin{textsample}{POS Dim 1 – 2014 – Score 18.00 – 2014/t001794.txt}  \label{ex:f1_pos_047}
Every day we face issues like climate change or \textbf{the} safety of vaccines where we have to answer questions whose answers rely heavily on scientific information.
Scientists tell us that \textbf{the} world is warming.
Scientists tell us that vaccines are safe.
But how do we know if they are right?
Why should be believe \textbf{the} science? \textbf{The} fact is, many of us actually don ' t believe \textbf{the} science.
Public opinion polls consistently show that significant proportions of \textbf{the} American people don ' t believe \textbf{the} climate is warming due to human activities, don ' t think that there is \textbf{evolution} by natural \textbf{selection}, and aren ' t persuaded by \textbf{the} safety of vaccines.
So why should we believe \textbf{the} science?
Well, scientists don ' t like talking about science as a matter of belief.
In fact, they would contrast science with faith, and they would say belief is \textbf{the} domain of faith.
And faith is a separate thing apart and distinct from science.
Indeed they would say religion is based on faith or maybe \textbf{the} calculus of Pascal ' s wager.
Blaise Pascal was a 17th-century mathematician who tried to bring scientific reasoning to \textbf{the} question of whether or not he should believe in God, and his wager went like this : Well, if God doesn ' t exist but I decide to believe in him nothing much is really lost.
Maybe a few hours on Sunday.
But if he does exist and I don ' t believe in him, then I ' m in deep trouble.
And so Pascal said, we ' d better believe in God.
Or as one of my college professors said,"He clutched for \textbf{the} handrail of faith."He made that leap of faith leaving science and rationalism behind.
Now \textbf{the} fact is though, for most of us, most scientific claims are a leap of faith.
We can ' t really judge scientific claims for ourselves in most cases.
And indeed this is actually true for most scientists as well outside of their own specialties.
So if you think about it, a geologist can ' t tell you whether a vaccine is safe.
Most chemists are not experts in \textbf{evolutionary} theory.
A physicist cannot tell you, despite \textbf{the} claims of some of them, whether or not tobacco causes cancer.
So, if even scientists themselves have to make a leap of faith outside their own fields, then why do they accept \textbf{the} claims of other scientists?
Why do they believe each other ' s claims?
And should we believe those claims?
So what I ' d like to argue is yes, we should, but not for \textbf{the} reason that most of us think.
Most of us were taught in school that \textbf{the} reason we should believe in science is because of \textbf{the} scientific method.
We were taught that scientists follow a method and that this method guarantees \textbf{the} truth of their claims. \textbf{The} method that most of us were taught in school, we can call it \textbf{the} \textbf{textbook} method, is \textbf{the} hypothetical deductive method.
According to \textbf{the} standard model, \textbf{the} \textbf{textbook} model, scientists develop hypotheses, they deduce \textbf{the} consequences of those hypotheses, and then they go out into \textbf{the} world and they say,"Okay, well are those consequences true?"Can we observe them taking place in \textbf{the} natural world?
And if they are true, then \textbf{the} scientists say,"Great, we know \textbf{the} hypothesis is correct."So there are many \textbf{famous} examples in \textbf{the} history of science of scientists doing exactly this.
One of \textbf{the} most \textbf{famous} examples comes from \textbf{the} work of Albert Einstein.
When Einstein developed \textbf{the} theory of general relativity, one of \textbf{the} consequences of his theory was that space-time wasn ' t just an empty void but that it actually had a fabric.
And that that fabric was bent in \textbf{the} presence of massive \textbf{objects} like \textbf{the} sun.
So if this theory were true then it meant that light as it passed \textbf{the} sun should actually be bent around it.
That was a pretty startling prediction and it took a few years before scientists were able to test it but they did test it in 1919, and lo and behold it turned out to be true.
Starlight actually does bend as it travels around \textbf{the} sun.
This was a huge confirmation of \textbf{the} theory.
It was considered proof of \textbf{the} truth of this radical new idea, and it was written up in many newspapers around \textbf{the} globe.
Now, sometimes this theory or this model is referred to as \textbf{the} deductive- nomological model, mainly because academics like to make things complicated.
But also because in \textbf{the} ideal case, it ' s about laws.
So nomological means having to do with laws.
And in \textbf{the} ideal case, \textbf{the} hypothesis isn ' t just an idea : ideally, it is a law of nature.
Why does it matter that it is a law of nature?
Because if it is a law, it can ' t be broken.
If it ' s a law then it will always be true in all times and all places no matter what \textbf{the} circumstances are.
And all of you know of at least one example of a \textbf{famous} law : Einstein ' s \textbf{famous} equation, E=MC2, which tells us what \textbf{the} relationship is between energy and mass.
And that relationship is true no matter what.
Now, it turns out, though, that there are several problems with this model. \textbf{The} main problem is that it ' s wrong.
It ' s just not true.
And I ' m going to talk about three reasons why it ' s wrong.
So \textbf{the} first reason is a logical reason.
It ' s \textbf{the} problem of \textbf{the} fallacy of affirming \textbf{the} consequent.
So that ' s another fancy, academic way of saying that false theories can make true predictions.
So just because \textbf{the} prediction comes true doesn ' t actually logically prove that \textbf{the} theory is correct.
And I have a good example of that too, again from \textbf{the} history of science.
This is a picture of \textbf{the} Ptolemaic \textbf{universe} with \textbf{the} Earth at \textbf{the} center of \textbf{the} \textbf{universe} and \textbf{the} sun and \textbf{the} \textbf{planets} going around it. \textbf{The} Ptolemaic model was believed by many very smart people for many centuries.
Well, why?
Well \textbf{the} answer is because it made lots of predictions that came true. \textbf{The} Ptolemaic system enabled astronomers to make accurate predictions of \textbf{the} motions of \textbf{the} \textbf{planet}, in fact more accurate predictions at first than \textbf{the} Copernican theory which we now would say is true.
So that ' s one problem with \textbf{the} \textbf{textbook} model.
A second problem is a practical problem, and it ' s \textbf{the} problem of auxiliary hypotheses.
Auxiliary hypotheses are assumptions that scientists are making that they may or may not even be aware that they ' re making.
So an important example of this comes from \textbf{the} Copernican model, which ultimately replaced \textbf{the} Ptolemaic system.
So when Nicolaus Copernicus said, actually \textbf{the} Earth is not \textbf{the} center of \textbf{the} \textbf{universe}, \textbf{the} sun is \textbf{the} center of \textbf{the} solar system, \textbf{the} Earth moves around \textbf{the} sun.
Scientists said, well \textbf{okay}, Nicolaus, if that ' s true we ought to be able to \textbf{detect} \textbf{the} motion of \textbf{the} Earth around \textbf{the} sun.
And so this \textbf{slide} here illustrates a concept known as stellar parallax.
And astronomers said, if \textbf{the} Earth is moving and we look at a prominent star, let ' s say, Sirius — well I know I ' m in Manhattan so you guys can ' t see \textbf{the} stars, but imagine you ' re out in \textbf{the} country, imagine you chose that rural life — and we look at a star in December, we see that star against \textbf{the} backdrop of distant stars.
If we now make \textbf{the} same observation six months later when \textbf{the} Earth has moved to this position in June, we look at that same star and we see it against a different backdrop.
That difference, that angular difference, is \textbf{the} stellar parallax.
So this is a prediction that \textbf{the} Copernican model makes.
Astronomers looked for \textbf{the} stellar parallax and they found nothing, nothing at all.
And many people argued that this proved that \textbf{the} Copernican model was false.
So what happened?
Well, in hindsight we can say that astronomers were making two auxiliary hypotheses, both of which we would now say were incorrect. \textbf{The} first was an assumption about \textbf{the} size of \textbf{the} Earth ' s orbit.
Astronomers were assuming that \textbf{the} Earth ' s orbit was large \textbf{relative} to \textbf{the} distance to \textbf{the} stars.
Today we would draw \textbf{the} picture more like this, this comes from NASA, and you see \textbf{the} Earth ' s orbit is actually quite small.
In fact, it ' s actually much smaller even than shown here. \textbf{The} stellar parallax therefore, is very small and actually very hard to \textbf{detect}.
And that leads to \textbf{the} second reason why \textbf{the} prediction didn ' t work, because scientists were also assuming that \textbf{the} telescopes they had were sensitive enough to \textbf{detect} \textbf{the} parallax.
And that turned out not to be true.
It wasn ' t until \textbf{the} 19th century that scientists were able to \textbf{detect} \textbf{the} stellar parallax.
So, there ' s a third problem as well. \textbf{The} third problem is simply a factual problem, that a lot of science doesn ' t fit \textbf{the} \textbf{textbook} model.
A lot of science isn ' t deductive at all, it ' s actually inductive.
And by that we mean that scientists don ' t necessarily start with theories and hypotheses, often they just start with observations of stuff going on in \textbf{the} world.
And \textbf{the} most \textbf{famous} example of that is one of \textbf{the} most \textbf{famous} scientists who ever lived, Charles Darwin.
When Darwin went out as a young man on \textbf{the} voyage of \textbf{the} Beagle, he didn ' t have a hypothesis, he didn ' t have a theory.
He just knew that he wanted to have a career as a scientist and he started to collect data.
Mainly he knew that he hated medicine because \textbf{the} sight of \textbf{blood} made him sick so he had to have an alternative career path.
So he started collecting data.
And he collected many things, including his \textbf{famous} finches.
When he collected these finches, he threw them in a bag and he had no idea what they meant.
Many years later back in London, Darwin looked at his data again and began to develop an explanation, and that explanation was \textbf{the} theory of natural \textbf{selection}.
Besides inductive science, scientists also often participate in modeling.
One of \textbf{the} things scientists want to do in life is to explain \textbf{the} causes of things.
And how do we do that?
Well, one way you can do it is to build a model that tests an idea.
So this is a picture of Henry Cadell, who was a Scottish geologist in \textbf{the} 19th century.
You can tell he ' s Scottish because he ' s wearing a deerstalker cap and Wellington boots.
And Cadell wanted to answer \textbf{the} question, how are mountains formed?
And one of \textbf{the} things he had observed is that if you look at mountains like \textbf{the} Appalachians, you often find that \textbf{the} \textbf{rocks} in them are folded, and they ' re folded in a particular way, which suggested to him that they were actually being compressed from \textbf{the} side.
And this idea would later play a major role in discussions of continental drift.
So he built this model, this crazy contraption with levers and wood, and here ' s his wheelbarrow, buckets, a big sledgehammer.
I don ' t know why he ' s got \textbf{the} Wellington boots.
Maybe it ' s going to rain.
And he created this physical model in order to demonstrate that you could, in fact, create patterns in \textbf{rocks}, or at least, in this case, in mud, that looked a lot like mountains if you compressed them from \textbf{the} side.
So it was an argument about \textbf{the} cause of mountains.
Nowadays, most scientists prefer to work inside, so they don ' t build physical models so much as to make \textbf{computer} simulations.
But a \textbf{computer} simulation is a kind of a model.
It ' s a model that ' s made with mathematics, and like \textbf{the} physical models of \textbf{the} 19th century, it ' s very important for thinking about causes.
So one of \textbf{the} big questions to do with climate change, we have tremendous amounts of evidence that \textbf{the} Earth is warming up.
This \textbf{slide} here, \textbf{the} black line shows \textbf{the} measurements that scientists have taken for \textbf{the} last 150 years showing that \textbf{the} Earth ' s temperature has steadily increased, and you can see in particular that in \textbf{the} last 50 years there ' s been this dramatic increase of nearly one degree centigrade, or almost two degrees Fahrenheit.
So what, though, is driving that change?
How can we know what ' s causing \textbf{the} observed warming?
Well, scientists can model it using a \textbf{computer} simulation.
So this diagram illustrates a \textbf{computer} simulation that has looked at all \textbf{the} different factors that we know can influence \textbf{the} Earth ' s climate, so sulfate particles from air pollution, volcanic dust from volcanic eruptions, changes in solar radiation, and, of course, \textbf{greenhouse} gases.
And they asked \textbf{the} question, what set of variables put into a model will reproduce what we actually see in real life?
So here is \textbf{the} real life in black.
Here ' s \textbf{the} model in this light gray, and \textbf{the} answer is a model that includes, it ' s \textbf{the} answer E on that SAT, all of \textbf{the} above. \textbf{The} only way you can reproduce \textbf{the} observed temperature measurements is with all of these things put together, including \textbf{greenhouse} gases, and in particular you can see that \textbf{the} increase in \textbf{greenhouse} gases tracks this very dramatic increase in temperature over \textbf{the} last 50 years.
And so this is why climate scientists say it ' s not just that we know that climate change is happening, we know that \textbf{greenhouse} gases are a major part of \textbf{the} reason why.
So now because there all these different things that scientists do, \textbf{the} philosopher Paul Feyerabend famously said,"\textbf{The} only principle in science that doesn ' t inhibit progress is : anything goes."Now this quotation has often been taken out of context, because Feyerabend was not actually saying that in science anything goes.
What he was saying was, actually \textbf{the} full quotation is,"If you press me to say what is \textbf{the} method of science, I would have to say : anything goes."What he was trying to say is that scientists do a lot of different things.
Scientists are creative.
But then this pushes \textbf{the} question back : If scientists don ' t use a single method, then how do they decide what ' s right and what ' s wrong?
And who judges?
And \textbf{the} answer is, scientists judge, and they judge by judging evidence.
Scientists collect evidence in many different ways, but however they collect it, they have to subject it to scrutiny.
And this led \textbf{the} sociologist Robert Merton to focus on this question of how scientists scrutinize data and evidence, and he said they do it in a way he called"organized skepticism."And by that he meant it ' s organized because they do it collectively, they do it as a group, and skepticism, because they do it from a position of distrust.
That is to say, \textbf{the} burden of proof is on \textbf{the} person with a novel claim.
And in this sense, science is intrinsically conservative.
It ' s quite hard to persuade \textbf{the} scientific community to say,"Yes, we know something, this is true."So despite \textbf{the} popularity of \textbf{the} concept of paradigm shifts, what we find is that actually, really major changes in scientific thinking are relatively rare in \textbf{the} history of science.
So finally that brings us to one more idea : If scientists judge evidence collectively, this has led historians to focus on \textbf{the} question of consensus, and to say that at \textbf{the} end of \textbf{the} day, what science is, what scientific knowledge is, is \textbf{the} consensus of \textbf{the} scientific experts who through this process of organized scrutiny, collective scrutiny, have judged \textbf{the} evidence and come to a conclusion about it, either yea or nay.
So we can think of scientific knowledge as a consensus of experts.
We can also think of science as being a kind of a jury, except it ' s a very special kind of jury.
It ' s not a jury of your peers, it ' s a jury of \textbf{geeks}.
It ' s a jury of men and women with Ph.
D. s, and unlike a conventional jury, which has only two choices, guilty or not guilty, \textbf{the} scientific jury actually has a number of choices.
Scientists can say yes, something ' s true.
Scientists can say no, it ' s false.
Or, they can say, well it might be true but we need to work more and collect more evidence.
Or, they can say it might be true, but we don ' t know how to answer \textbf{the} question and we ' re going to put it aside and maybe we ' ll come back to it later.
That ' s what scientists call"intractable."But this leads us to one final problem : If science is what scientists say it is, then isn ' t that just an appeal to authority?
And weren ' t we all taught in school that \textbf{the} appeal to authority is a logical fallacy?
Well, here ' s \textbf{the} paradox of modern science, \textbf{the} paradox of \textbf{the} conclusion I think historians and philosophers and sociologists have come to, that actually science is \textbf{the} appeal to authority, but it ' s not \textbf{the} authority of \textbf{the} individual, no matter how smart that individual is, like Plato or Socrates or Einstein.
It ' s \textbf{the} authority of \textbf{the} collective community.
You can think of it is a kind of wisdom of \textbf{the} crowd, but a very special kind of crowd.
Science does appeal to authority, but it ' s not based on any individual, no matter how smart that individual may be.
It ' s based on \textbf{the} collective wisdom, \textbf{the} collective knowledge, \textbf{the} collective work, of all of \textbf{the} scientists who have worked on a particular problem.
Scientists have a kind of culture of collective distrust, this"show me"culture, illustrated by this nice woman here showing her colleagues her evidence.
Of course, these people don ' t really look like scientists, because they ' re much too happy. \textbf{Okay}, so that brings me to my final point.
Most of us get up in \textbf{the} morning.
Most of us trust our cars.
Well, see, now I ' m thinking, I ' m in Manhattan, this is a bad analogy, but most Americans who don ' t live in Manhattan get up in \textbf{the} morning and get in their cars and turn on that ignition, and their cars work, and they work incredibly well. \textbf{The} modern automobile hardly ever breaks down.
So why is that?
Why do cars work so well?
It ' s not because of \textbf{the} genius of Henry Ford or Karl Benz or even Elon Musk.
It ' s because \textbf{the} modern automobile is \textbf{the} product of more than 100 years of work by hundreds and thousands and tens of thousands of people. \textbf{The} modern automobile is \textbf{the} product of \textbf{the} collected work and wisdom and experience of every man and woman who has ever worked on a car, and \textbf{the} reliability of \textbf{the} technology is \textbf{the} result of that accumulated effort.
We benefit not just from \textbf{the} genius of Benz and Ford and Musk but from \textbf{the} collective intelligence and hard work of all of \textbf{the} people who have worked on \textbf{the} modern car.
And \textbf{the} same is true of science, only science is even older.
Our basis for trust in science is actually \textbf{the} same as our basis in trust in technology, and \textbf{the} same as our basis for trust in anything, namely, experience.
But it shouldn ' t be blind trust any more than we would have blind trust in anything.
Our trust in science, like science itself, should be based on evidence, and that means that scientists have to become better communicators.
They have to explain to us not just what they know but how they know it, and it means that we have to become better listeners.
Thank you very much.

% matched lemmas: blood, computer, detect, evolution, evolutionary, famous, geek, greenhouse, object, okay, planet, relative, rock, selection, slide, textbook, the, universe
\end{textsample}
